\scanpage{315}\mspage{39}
$p,q\in\mathcal F_W$ et $f\in W$, on a $f\in W^{\unsure{u|v}}$ (ou,
ce qui revient au même, \deleted{\unsure{cf $\overline{\underline{\mathrm{VI}}}$}} comme $X,Y$ sont
$W$-lisses sur $S$, $f_s\in W$ pour tt $s\in\mathrm{Ob}\, S$).
% slanted
\margin{\uline{NB} La \unsure{dém.} de $\mathrm B(\mathcal F_W)\Longrightarrow\mathrm B(\mathrm{Fib}f_W)$ montre en fait qu'il suffit de postuler b) lorsque $f$ est cartésien, et $=$ lorsque $X,Y$ sont scindées sur $S$ et $f$ compatible avec les scindages.}

c) \deleted{\unsure{Pour tte flèche}} On a
\[
\mathcal F_W\cap W\subset W^u\,,
\]
en d'autres termes, pour toute flèche
$f:X\unsure{\mathrel{\text{---}}|\longrightarrow} S$ dans $\mathrm{Cat}$ telle que $f\in\mathcal F_W$
(on pourrait appeler ces flèches des Cat-$W$-fibrations, étant entendu que \unsure{cela|elles} sont \inserted{aussi} des
Cat-fibrations, mais non \deleted{\unsure{f}}nécessairement
des $W$-fibrations), si $f\in W$, alors
$f\in W_S^u$ (ou, ce qui revient au $=$, les fibres
$X_s$ de $f$ ($s\in\mathrm{Ob}\,S$) sont $W$-asphériques).

% slanted
\margin{rajouté après coup $\rightsquigarrow$}
\inserted{F) $\mathrm{Cat}$ a suffisamment de $W$-fibrations, i.e. $\mathrm{Cat}$
munie de $(W,\mathrm{Fib}_W)$ est une cat. à cofibrations (propres, et
$=$ de \deleted{\unsure{Br}}\unsure{Brown}, i.e. $X\longrightarrow e$ est dans $\mathrm{Fib}_W$ pour tt $X$ dans $\mathrm{Cat}$).}

\unsure{Il} est clair que $\unsure{b|c}$ est \unsure{un} cas particulier
de $\unsure{c|b}$ obtenu en prenant $Y=S$, $q=\mathrm{id}_S$,
donc $b\Longrightarrow c$. L'implication $a\Longrightarrow b$ est
moins évidente, elle est établie d'abord dans
$\overline{\underline{\mathrm{XII}}}$, et nous y reviendrons ci-dessous \inserted{(4.4.2 ci-dessous)}.
Le reste du présent \S\ est consacré à \deleted{\unsure{donner}} \inserted{étudier}
diverses variantes des conditions a) b) c)
précédentes, et les implications \deleted{entre} celles-ci.
% archive sheet number: 314
